% !TEX TS-program = xelatex
% !TEX encoding = UTF-8 Unicode

\documentclass{resume}
\usepackage{zh_CN-Adobefonts_external}
\usepackage{linespacing_fix}
\usepackage{cite}
\usepackage{geometry}
\geometry{a4paper,left=0.5cm,right=0.5cm,top=0.5cm,bottom=0.5cm}

\begin{document}
\pagenumbering{gobble}

\name{张宇航 (Yuhang Zhang)}
\basicInfo{
	\email{yuhang.zhang@example.com} \textperiodcentered\ 
	\phone{+86 138-1234-5678} \textperiodcentered\ 
	\faCalendar{2001年5月18日} \textperiodcentered\ 
	\faGroup {男} \textperiodcentered\ 
	\faGithub{\href{https://github.com/zhangyuhang-dev}{zhangyuhang-dev}}
}

%% 教育背景
\section{教育背景}

\datedsubsection{\textbf{爱丁堡大学 (University of Edinburgh)} \ 电子学 \ 硕士学位}{2024年9月 -- 2025年12月}
\begin{itemize}
    \item \textit{GPA：}\ 3.85/4.00\ \textit{排名：}\ 前15\%\ \textit{相关课程：}\ 实时系统设计、自主机器人系统、统计信号处理、模拟集成电路设计、纳米电子器件
\end{itemize}

\datedsubsection{\textbf{南京邮电大学} \ 电子信息工程 \ 学士学位}{2020年9月 -- 2024年6月}
\begin{itemize}
    \item \textit{GPA：}\ 84.50/100\ \textit{排名：}\ 前30\%\ \textit{相关课程：}\ 信号与系统、微机原理与接口技术、计算机网络、模拟与数字电子技术、电磁场与电磁波
\end{itemize}

% =============================================================================
% 技能与兴趣
% =============================================================================
\section{\textbf{技能与兴趣}}
\begin{onehalfspacing}
\begin{itemize}
    \item \textbf{编程语言与工具：} 熟悉 C/C++ (多线程开发)；熟练使用 Python 进行数据处理与脚本自动化；熟悉 Verilog/VHDL ；掌握 Shell/Bash 脚本、Makefile 构建、Git 版本控制及 GDB 调试。
    \item \textbf{嵌入式 Linux 开发：} 了解 Linux 内核，具备 BSP 开发能力 (U-Boot, Device Tree, Rootfs)；有驱动开发 (字符设备) 开发经验；熟悉 Buildroot 构建环境。
    \item \textbf{硬件与 FPGA 设计：} 熟练使用 AD 进行 PCB 设计；熟悉 FPGA 异构开发流程；精通 MCU 及 FreeRTOS 系统应用。
    \item \textbf{人工智能与边缘计算：} 熟悉 PyTorch 深度学习框架。
\end{itemize}
\end{onehalfspacing}

%% 科研经历
\section{项目经历}

\datedsubsection{\textbf{基于多模态传感器融合与轻量化时序网络的智能机械故障预测系统} \ 毕业设计 (独立研究员)}{2025年6月 -- 2026年8月}
\begin{onehalfspacing}
    \textbf{项目简介：} 针对工业物联网中旋转机械的状态监测需求，设计并开发了一套基于深度学习的多通道振动与温度信号融合诊断系统，实现轴承早期微弱故障的精准预测与健康评估。\\
    \textbf{主要职责：} 主导多模态数据采集与轻量化时序网络模型设计，实现边缘网关故障分析开发。\\
    \textbf{亮点：}
    \begin{itemize}
        \item \textbf{数据采集与去噪：} 负责多通道加速度与温度信号同步采集、小波包去噪及特征标准化。
        \item \textbf{网络模型设计：} 设计结合时间卷积网络 (TCN) 与自注意力机制的轻量化多源时序信号融合模型。
        \item \textbf{模型量化与部署：} 在 PyTorch 框架下完成网络训练，应用剪枝与 8-bit 量化技术部署于边缘网关。
        \item \textbf{分类准确度：} 在公开及实测数据集上实现了 10 种典型机械状态分类，平均准确率达 94.2\%。
    \end{itemize}
\end{onehalfspacing}

\datedsubsection{\textbf{基于 RT-Thread 与多节点低功耗无线网络的智能微气候监测系统} \ 系统架构师与底层开发}{2025年3月 -- 2026年6月}
\begin{onehalfspacing}
    \textbf{项目简介：} 研发面向智慧农业微气候环境的高精度、低功耗分布式监测网络，实现多节点空气温湿度、光照、土壤湿度等数据的自组网采集与云端联动。\\
    \textbf{主要职责：} 负责多路传感器驱动编写、Sub-1G/LoRa 自组网协议设计及网关系统移植开发。\\
    \textbf{亮点：}
    \begin{itemize}
        \item \textbf{超低功耗硬件：} 基于低功耗 SoC 芯片自主设计多层 PCB，优化电源管理，使休眠电流降至微安级。
        \item \textbf{RTOS 驱动开发：} 移植 RT-Thread 实时操作系统，编写多通道 ADC、I2C 传感器及电源管理底层驱动。
        \item \textbf{自组网协议设计：} 编写基于 Sub-1G/LoRa 的无线协议，设计抗干扰时分多址 (TDMA) 数据上传机制。
        \item \textbf{云端系统集成：} 搭建 MQTT 网关，对接云端物联网平台，设计基于微信小程序的可视化系统。
    \end{itemize}
\end{onehalfspacing}

%% 所获荣誉
\section{所获荣誉}
\datedsubsection{\textbf{2022年9月--2023年9月}}{\ 全国大学生英语竞赛 (国家级)}
\datedsubsection{\textbf{2021年9月--2022年9月}}{\ 电子设计竞赛（校级） (校级)}

\end{document}